% 05-materiales-metodos.tex — Fuente: docs/informe-final/05-materiales-metodos.md

\chapter{Materiales y métodos}
\label{cap:materiales-metodos}

\section{Enfoque metodológico global}

El proyecto sigue la metodología Design Science Research (\gls{dsr}) de
Peffers et al.~\cite{PEFFERS2007} en sus seis actividades canónicas,
instanciadas así:

\begin{table}[H]
\centering
\caption{Instanciación de las seis actividades DSR de Peffers et al.}
\label{tab:dsr}
\begin{tabular}{@{}p{4.6cm}p{9.4cm}@{}}
\toprule
\textbf{Actividad DSR} & \textbf{Instanciación en Artisync} \\
\midrule
1. Identificación del problema y motivación & Capítulo~\ref{cap:introduccion}
  --- fragmentación de la economía creativa digital, ausencia de
  garantías de pago/identidad \\
2. Definición de objetivos de una solución & Objetivos específicos del
  Capítulo~\ref{cap:introduccion}~\S1.3, derivados de las RQ \\
3. Diseño y desarrollo & Capítulos~\ref{cap:diseno-arquitectura}
  (arquitectura) y~\ref{cap:implementacion} (implementación); cuatro
  entregas incrementales (1A$\to$1B$\to$3$\to$Final) \\
4. Demostración & Sistema desplegado y operativo, con usuario demo
  documentado (\texttt{README.md}) \\
5. Evaluación & Capítulo~\ref{cap:evaluacion-resultados} --- evaluación
  empírica multidimensional (rendimiento, seguridad, usabilidad,
  cobertura, calidad web) \\
6. Comunicación & Este documento, el repositorio público con DOI Zenodo, y
  la defensa oral del PFC \\
\bottomrule
\end{tabular}
\end{table}

El propio marco DSR de Peffers et al. se apoya en los siete lineamientos
de rigor de Hevner, March, Park y Ram~\cite{HEVNER2004} para investigación
en sistemas de información (relevancia del problema, rigor de diseño,
evaluación del artefacto, contribución a la investigación, rigor de
investigación, diseño como proceso de búsqueda, comunicación de
resultados), que este proyecto usa como criterio implícito de calidad al
decidir qué evidencia archivar en \texttt{docs/mediciones/}.

Dado que el trabajo aporta evaluación empírica de rendimiento y de
usabilidad de un artefacto único (no un experimento controlado con grupos
de comparación aleatorizados ni un estudio de caso sobre una organización
externa), se acoge además el estándar empírico \emph{Engineering
Research} de Ralph et al.~\cite{RALPH2021} para el reporte metodológico
--- ver checklist completo en \texttt{docs/checklists/ralph-2021-checklist.md}.
Se consultaron adicionalmente las guías de reporte de estudios de caso de
Runeson y Höst~\cite{RUNESON2009} como referencia de rigor metodológico
para documentar el protocolo experimental (\S5.5), aunque Artisync no se
reporta formalmente como un estudio de caso en el sentido estricto de esa
guía (no hay una organización externa como unidad de análisis).

\section{Instrumentos de medición}

\begin{table}[H]
\centering
\caption{Instrumentos de medición empírica utilizados}
\label{tab:instrumentos}
\begin{tabular}{@{}p{2.6cm}p{2.6cm}p{3.2cm}p{5.0cm}@{}}
\toprule
\textbf{Instrumento} & \textbf{Versión exacta} & \textbf{Qué mide} & \textbf{Configuración} \\
\midrule
SUS de Brooke~\cite{BROOKE1996,BANGOR2009} & 10 ítems originales, sin
  modificación & Usabilidad percibida & Formulario Google Forms, escala
  1--5, ver \texttt{docs/\allowbreak mediciones/\allowbreak sus/\allowbreak
  instrucciones-\allowbreak formulario.md} \\
k6 & v2.1.0 (go1.26.4) & Rendimiento bajo carga & 50 VUs, 30 s,
  \texttt{k6/catalogo-load.js} (ver \S5.5) \\
Lighthouse / \texttt{@lhci/cli} & Lighthouse 12.6.1 / \texttt{@lhci/cli}
  0.15.1 & Calidad web (Performance, Accessibility, Best Practices, SEO) &
  \texttt{lighthouserc.\allowbreak mobile.json} / \texttt{lighthouserc.\allowbreak
  desktop.json}, 3 corridas por perfil \\
JaCoCo~\cite{ZHU1997,INOZEMTSEVA2014} & 0.8.13 (jacoco-maven-plugin) &
  Cobertura de código (líneas, ramas, complejidad) & \texttt{pom.xml},
  umbral declarado $\geq$70\,\% \\
OWASP ZAP~\cite{OWASP2021} & \texttt{ghcr.io/\allowbreak zaproxy/\allowbreak
  zaproxy:stable} (\texttt{zap-baseline.py}) & Escaneo dinámico de
  seguridad & Modo baseline contra \texttt{http://\allowbreak localhost:4200} \\
SpotBugs + find-sec-bugs~\cite{HALFOND2005,GOULD2004} & 4.8.6.6 + 1.13.0 &
  Análisis estático de inyección SQL & Regla
  \texttt{SQL\_NONCONSTANT\_STRING\_PASSED\_TO\_EXECUTE} y variantes \\
\bottomrule
\end{tabular}
\end{table}

Versiones exactas de todo el entorno (Docker, JDK, Node, Angular CLI) en
\texttt{docs/entorno/versions.txt}.

\section{Enfoque de métricas: Goal-Question-Metric (GQM)}

\textbf{Ejemplo de GQM completo aplicado al bloque de rendimiento}, siguiendo
el enfoque original de Basili, Caldiera y Rombach~\cite{BASILI1994}:

\begin{itemize}
  \item \textbf{Meta (Goal):} Analizar el rendimiento del endpoint de
    catálogo con el propósito de caracterizarlo desde el punto de vista
    de la latencia percibida por el usuario, en el contexto del sistema
    Artisync bajo carga concurrente moderada.
  \item \textbf{Preguntas (Questions):}
    \begin{itemize}
      \item Q1: ¿Cuál es la latencia p95 del endpoint bajo 50 usuarios
        virtuales concurrentes?
      \item Q2: ¿Se mantiene esa latencia dentro del umbral declarado
        (200~ms con caché caliente, 500~ms con caché frío)?
      \item Q3: ¿Existen errores de servidor (HTTP $\geq$500) bajo esa
        carga?
    \end{itemize}
  \item \textbf{Métricas (Metrics):}
    \begin{itemize}
      \item M1: percentil 95 de \texttt{http\_req\_duration} (ms), por
        escenario.
      \item M2: tasa de \texttt{http\_req\_failed} con código $\geq$500
        (\%).
      \item M3: throughput agregado (\texttt{http\_reqs}/s).
    \end{itemize}
\end{itemize}

El mismo patrón GQM (meta declarada en el \texttt{REPORTE-*.md} de cada
bloque, preguntas implícitas en los umbrales del SRS, métricas reportadas
con estadística descriptiva) se repite para seguridad, usabilidad,
cobertura y calidad web --- ver \texttt{docs/mediciones/<bloque>/REPORTE-*.md}.

\section{Población, muestreo y participantes}

\textbf{Bloque de usabilidad (SUS):} población objetivo --- usuarios
potenciales de una plataforma de intermediación de servicios digitales
(perfil generalista, no exclusivamente técnico). Muestra: \textbf{N=16
participantes externos al equipo de desarrollo}, reclutados por
conveniencia (red de contactos del equipo), rango de edad 21--62 años
(M=34.6), 50\,\% femenino / 50\,\% masculino, experiencia web 50\,\% alta
/ 31.3\,\% media / 18.8\,\% baja, dispositivo 37.5\,\% móvil / 25\,\%
laptop / 25\,\% desktop / 12.5\,\% tablet
(\texttt{docs/mediciones/sus/perfil-participantes.csv}).

\textbf{Limitación de muestreo declarada honestamente} (siguiendo Baltes
y Ralph~\cite{BALTES2022} sobre reporte de muestreo en ingeniería de
software empírica): el muestreo por conveniencia no garantiza
representatividad de la población de usuarios reales de la plataforma
(creadores y clientes de economía creativa específicamente), y
\textbf{no existe una justificación de tamaño de muestra basada en poder
estadístico} --- N=16 se decidió por disponibilidad de participantes
dentro del plazo académico, no por un cálculo a priori. Esta limitación
se retoma en el Capítulo~\ref{cap:validez} (Amenazas a la validez
externa).

\textbf{Bloques automatizados (k6, Lighthouse, JaCoCo, OWASP/ZAP):} no
aplican muestreo de sujetos humanos; el ``tamaño de muestra'' es el
número de corridas independientes (3 para Lighthouse por perfil, 3+3 para
k6 por escenario caliente/frío, 1 ejecución exhaustiva para JaCoCo sobre
la suite de 522 pruebas).

\section{Protocolo experimental replicable}

Cada bloque de evaluación tiene un comando reproducible documentado en
\texttt{Makefile} y su reporte:

\begin{itemize}
  \item \textbf{Rendimiento:} \texttt{make bench} (requiere \texttt{make
    up} y k6 instalado) --- ejecuta \texttt{k6/catalogo-load.js} contra
    \texttt{GET /api/v1/catalogo?page=0\&size=20}, 50~VUs/30s. Ver
    \texttt{docs/mediciones/perf/REPORTE-PERF.md} para el protocolo
    detallado de los escenarios caliente y frío.
  \item \textbf{Seguridad:} \texttt{make audit} (análisis estático) +
    \texttt{make audit-zap} (escaneo dinámico) --- ver
    \texttt{docs/mediciones/sec/REPORTE-SEC.md}.
  \item \textbf{Usabilidad:} protocolo de tarea guiada
    (registro/login $\to$ explorar catálogo $\to$ crear pedido $\to$
    revisar seguimiento $\to$ cerrar sesión) descrito en
    \texttt{docs/mediciones/sus/instrucciones-formulario.md}, seguido del
    cuestionario SUS; análisis con \texttt{make sus}.
  \item \textbf{Cobertura:} \texttt{make test} (JUnit~5 + JaCoCo, reporte
    en \texttt{docs/mediciones/jacoco/}).
  \item \textbf{Calidad web:} \texttt{make lighthouse} --- ver detalle de
    la configuración del contenedor efímero en el propio \texttt{Makefile}
    (razones técnicas del aislamiento de red documentadas ahí).
\end{itemize}

\section{Análisis estadístico}

Estadísticos calculados en todas las mediciones cuantitativas:
\textbf{media, desviación típica e intervalo de confianza al 95\,\%},
preferidos sobre valores p aislados como forma primaria de reporte,
siguiendo la recomendación de la comunidad de ingeniería de software
empírica~\cite{WOHLIN2012,ARCURI2011}. Ver valores concretos por bloque en
el Capítulo~\ref{cap:evaluacion-resultados}.

\textbf{Brecha declarada honestamente:} la guía de la Entrega Final exige,
cuando se comparan dos condiciones (ej. caché frío vs. caliente en el
bloque de rendimiento), un test inferencial no paramétrico (Wilcoxon,
Mann-Whitney) acompañado de un tamaño de efecto (\emph{Cliff's delta}, r
de rangos) --- exactamente el procedimiento que Arcuri y
Briand~\cite{ARCURI2011} sistematizan para algoritmos y comparaciones no
deterministas en ingeniería de software. \textbf{Ese test no se ejecutó}
--- \texttt{docs/checklists/ralph-2021-checklist.md} y
\texttt{docs/checklists/fair-checklist.md} documentan esta ausencia sin
maquillarla. Además, la propia comparación caliente/frío tiene una
limitación metodológica de diseño documentada en
\texttt{docs/mediciones/perf/REPORTE-PERF.md}: el script de carga no
aísla un \emph{cache miss} real por iteración, por lo que un test
inferencial sobre los datos actuales no respondería a la pregunta que
pretende responder. Corregir el diseño del experimento (variar los
parámetros de consulta para forzar \emph{misses} reales) es prerrequisito
para que el test inferencial tenga sentido --- declarado como trabajo
futuro en el Capítulo~\ref{cap:trabajo-futuro}.

No se realizaron comparaciones múltiples que requieran corrección (Holm,
Benjamini-Hochberg) en esta entrega.

\section{Determinismo y semillas}

Ver el detalle completo en \texttt{docs/entorno/versions.txt}, sección
``Semillas aleatorias''. Resumen: \texttt{fn\_seleccionar\_ganadores\_sorteo}
usa \texttt{random()} de PostgreSQL sin semilla fija (no determinista por
diseño --- es un sorteo real); los scripts de análisis SUS no usan
muestreo aleatorio. No se identificaron otras dependencias de semillas
aleatorias en los scripts de medición a la fecha de este documento.
